% Options for packages loaded elsewhere
% Options for packages loaded elsewhere
\PassOptionsToPackage{unicode}{hyperref}
\PassOptionsToPackage{hyphens}{url}
\PassOptionsToPackage{dvipsnames,svgnames,x11names}{xcolor}
%
\documentclass[
  spanish,
]{llncs}
\usepackage{xcolor}
\usepackage{amsmath,amssymb}
\setcounter{secnumdepth}{-\maxdimen} % remove section numbering
\usepackage{iftex}
\ifPDFTeX
  \usepackage[T1]{fontenc}
  \usepackage[utf8]{inputenc}
  \usepackage{textcomp} % provide euro and other symbols
\else % if luatex or xetex
  \usepackage{unicode-math} % this also loads fontspec
  \defaultfontfeatures{Scale=MatchLowercase}
  \defaultfontfeatures[\rmfamily]{Ligatures=TeX,Scale=1}
\fi
\usepackage{lmodern}
\ifPDFTeX\else
  % xetex/luatex font selection
\fi
% Use upquote if available, for straight quotes in verbatim environments
\IfFileExists{upquote.sty}{\usepackage{upquote}}{}
\IfFileExists{microtype.sty}{% use microtype if available
  \usepackage[]{microtype}
  \UseMicrotypeSet[protrusion]{basicmath} % disable protrusion for tt fonts
}{}
\makeatletter
\@ifundefined{KOMAClassName}{% if non-KOMA class
  \IfFileExists{parskip.sty}{%
    \usepackage{parskip}
  }{% else
    \setlength{\parindent}{0pt}
    \setlength{\parskip}{6pt plus 2pt minus 1pt}}
}{% if KOMA class
  \KOMAoptions{parskip=half}}
\makeatother
% Make \paragraph and \subparagraph free-standing
\makeatletter
\ifx\paragraph\undefined\else
  \let\oldparagraph\paragraph
  \renewcommand{\paragraph}{
    \@ifstar
      \xxxParagraphStar
      \xxxParagraphNoStar
  }
  \newcommand{\xxxParagraphStar}[1]{\oldparagraph*{#1}\mbox{}}
  \newcommand{\xxxParagraphNoStar}[1]{\oldparagraph{#1}\mbox{}}
\fi
\ifx\subparagraph\undefined\else
  \let\oldsubparagraph\subparagraph
  \renewcommand{\subparagraph}{
    \@ifstar
      \xxxSubParagraphStar
      \xxxSubParagraphNoStar
  }
  \newcommand{\xxxSubParagraphStar}[1]{\oldsubparagraph*{#1}\mbox{}}
  \newcommand{\xxxSubParagraphNoStar}[1]{\oldsubparagraph{#1}\mbox{}}
\fi
\makeatother


\usepackage{longtable,booktabs,array}
\newcounter{none} % for unnumbered tables
\usepackage{calc} % for calculating minipage widths
% Correct order of tables after \paragraph or \subparagraph
\usepackage{etoolbox}
\makeatletter
\patchcmd\longtable{\par}{\if@noskipsec\mbox{}\fi\par}{}{}
\makeatother
% Allow footnotes in longtable head/foot
\IfFileExists{footnotehyper.sty}{\usepackage{footnotehyper}}{\usepackage{footnote}}
\makesavenoteenv{longtable}
\usepackage{graphicx}
\makeatletter
\newsavebox\pandoc@box
\newcommand*\pandocbounded[1]{% scales image to fit in text height/width
  \sbox\pandoc@box{#1}%
  \Gscale@div\@tempa{\textheight}{\dimexpr\ht\pandoc@box+\dp\pandoc@box\relax}%
  \Gscale@div\@tempb{\linewidth}{\wd\pandoc@box}%
  \ifdim\@tempb\p@<\@tempa\p@\let\@tempa\@tempb\fi% select the smaller of both
  \ifdim\@tempa\p@<\p@\scalebox{\@tempa}{\usebox\pandoc@box}%
  \else\usebox{\pandoc@box}%
  \fi%
}
% Set default figure placement to htbp
\def\fps@figure{htbp}
\makeatother



\ifLuaTeX
\usepackage[bidi=basic,shorthands=off]{babel}
\else
\usepackage[bidi=default,shorthands=off]{babel}
\fi
\ifLuaTeX
  \usepackage{selnolig} % disable illegal ligatures
\fi


\setlength{\emergencystretch}{3em} % prevent overfull lines

\providecommand{\tightlist}{%
  \setlength{\itemsep}{0pt}\setlength{\parskip}{0pt}}



 


\makeatletter
\@ifpackageloaded{caption}{}{\usepackage{caption}}
\AtBeginDocument{%
\ifdefined\contentsname
  \renewcommand*\contentsname{Tabla de contenidos}
\else
  \newcommand\contentsname{Tabla de contenidos}
\fi
\ifdefined\listfigurename
  \renewcommand*\listfigurename{Listado de Figuras}
\else
  \newcommand\listfigurename{Listado de Figuras}
\fi
\ifdefined\listtablename
  \renewcommand*\listtablename{Listado de Tablas}
\else
  \newcommand\listtablename{Listado de Tablas}
\fi
\ifdefined\figurename
  \renewcommand*\figurename{Figura}
\else
  \newcommand\figurename{Figura}
\fi
\ifdefined\tablename
  \renewcommand*\tablename{Tabla}
\else
  \newcommand\tablename{Tabla}
\fi
}
\@ifpackageloaded{float}{}{\usepackage{float}}
\floatstyle{ruled}
\@ifundefined{c@chapter}{\newfloat{codelisting}{h}{lop}}{\newfloat{codelisting}{h}{lop}[chapter]}
\floatname{codelisting}{Listado}
\newcommand*\listoflistings{\listof{codelisting}{Listado de Listados}}
\makeatother
\makeatletter
\makeatother
\makeatletter
\@ifpackageloaded{caption}{}{\usepackage{caption}}
\@ifpackageloaded{subcaption}{}{\usepackage{subcaption}}
\makeatother
\usepackage{bookmark}
\IfFileExists{xurl.sty}{\usepackage{xurl}}{} % add URL line breaks if available
\urlstyle{same}
\hypersetup{
  pdftitle={Tarea Individual DG\_S4},
  pdfauthor={Justin Guevara},
  pdflang={es},
  colorlinks=true,
  linkcolor={blue},
  filecolor={Maroon},
  citecolor={Blue},
  urlcolor={Blue},
  pdfcreator={LaTeX via pandoc}}


\title{Tarea Individual DG\_S4}
\usepackage{etoolbox}
\makeatletter
\providecommand{\subtitle}[1]{% add subtitle to \maketitle
  \apptocmd{\@title}{\par {\large #1 \par}}{}{}
}
\makeatother
\subtitle{Formulación y Búsqueda Ciega en Espacios de Estados}
\author{Justin Guevara}
\date{}
\begin{document}
\maketitle


\textbf{Enlace al repositorio de GitHub}:
\url{https://github.com/JUSTAILS/IA---JUSTIN-GUEVARA}

\subsection{Objetivos}\label{objetivos}

\subsubsection{General}\label{general}

\begin{itemize}
\tightlist
\item
  Formular un problema mediante un espacio de estados e implementar
  algoritmos de búsqueda ciega para resolverlo, comparando las
  estrategias BFS y DFS mediante métricas de exploración, tiempo de
  ejecución y longitud de la solución.
\end{itemize}

\subsubsection{Específicos}\label{especuxedficos}

\begin{itemize}
\tightlist
\item
  Definir los componentes formales del problema: estados, estado
  inicial, estado objetivo, acciones, función de transición y
  restricciones de validez.
\item
  Implementar los algoritmos BFS y DFS en Python para explorar el
  espacio de estados y encontrar una ruta hacia el objetivo.
\item
  Medir el comportamiento de ambos algoritmos mediante el número de
  nodos expandidos, los estados descubiertos, el tiempo de ejecución, la
  longitud de la solución y el factor de ramificación efectivo.
\item
  Comparar los resultados experimentales para analizar las
  características y diferencias de ambas estrategias de búsqueda ciega.
\end{itemize}

\subsection{Introducción}\label{introducciuxf3n}

La inteligencia artificial comprende métodos que permiten representar
problemas y encontrar soluciones mediante procedimientos
computacionales. Entre sus fundamentos se encuentra la formulación de
problemas como espacios de estados, en los cuales una configuración
inicial puede transformarse en otras configuraciones mediante acciones
definidas, hasta alcanzar un estado que satisfaga las condiciones del
objetivo.

La búsqueda ciega, también denominada búsqueda no informada, constituye
una familia de estrategias que exploran el espacio de estados sin
emplear información heurística para estimar la proximidad al objetivo.
Dentro de esta categoría, la búsqueda en anchura (Breadth-First Search /
BFS) y la búsqueda en profundidad (Depth-First Search / DFS) representan
dos métodos relevantes, diferenciados principalmente por el orden en que
procesan los estados pendientes de exploración.

Para estudiar estos conceptos, se utiliza el problema de los misioneros
y caníbales, que consiste en trasladar tres misioneros y tres caníbales
entre las dos orillas de un río mediante una embarcación con capacidad
para dos personas. La solución debe respetar las restricciones que
impiden que los caníbales superen en número a los misioneros en una
orilla donde estos últimos estén presentes.

Este problema permite establecer formalmente el espacio de estados,
definir las acciones posibles y descartar las transiciones que producen
configuraciones no válidas. A partir de esta formulación, se implementan
BFS y DFS para encontrar rutas hacia el estado objetivo y comparar
experimentalmente su comportamiento.

\subsection{Formulación del problema}\label{formulaciuxf3n-del-problema}

\subsubsection{Descripción del
problema}\label{descripciuxf3n-del-problema}

\begin{itemize}
\tightlist
\item
  El problema de los misioneros y caníbales consiste en trasladar tres
  misioneros y tres caníbales desde la orilla izquierda hasta la orilla
  derecha de un río. Para realizar los cruces se dispone de una
  embarcación con capacidad para una o dos personas.
\end{itemize}

La solución debe satisfacer las siguientes condiciones:

\begin{enumerate}
\def\labelenumi{\arabic{enumi}.}
\tightlist
\item
  La embarcación solo puede transportar uno o dos pasajeros por cruce.
\item
  La embarcación debe encontrarse en la orilla desde la cual se inicia
  cada movimiento.
\item
  En ninguna orilla donde haya misioneros puede haber más caníbales que
  misioneros.
\item
  La cantidad total de misioneros y caníbales debe conservarse durante
  todo el proceso.
\item
  El objetivo consiste en trasladar a las seis personas a la orilla
  derecha.
\end{enumerate}

Estas condiciones permiten establecer un espacio de estados finito en el
que los algoritmos de búsqueda pueden explorar las transiciones
permitidas hasta encontrar una solución.

\subsubsection{Representación de los
estados}\label{representaciuxf3n-de-los-estados}

\begin{itemize}
\tightlist
\item
  Cada estado se representa mediante una tupla de tres componentes:
\end{itemize}

\[
S=(M,C,B)
\]

Donde:

(M): número de misioneros en la orilla izquierda. (C): número de
caníbales en la orilla izquierda. (B): posición de la embarcación,
definida como:

B=1, si el barco se encuentra en la orilla izquierda, y B=0, si el barco
se encuentra en la orilla derecha.

Como existen tres misioneros y tres caníbales, las cantidades presentes
en la orilla derecha se calculan mediante:

\[
M_D=3-M
\] \[
C_D=3-C
\]

Esta representación permite describir la distribución de las personas y
la posición de la embarcación sin almacenar por separado la información
de ambas orillas.

\subsubsection{Estado inicial y estado
objetivo}\label{estado-inicial-y-estado-objetivo}

\begin{itemize}
\tightlist
\item
  El estado inicial representa la situación anterior al primer cruce,
  cuando todas las personas y la embarcación se encuentran en la orilla
  izquierda:
\end{itemize}

\[
S_0=(3,3,1)
\]

El estado objetivo se alcanza cuando todos los misioneros y caníbales
están en la orilla derecha y la embarcación también se encuentra en
ella:

\[
S_G=(0,0,0)
\]

Por tanto, el algoritmo debe encontrar una secuencia de transiciones
válidas que conecte (S\_0) con (S\_G).

\subsubsection{Acciones permitidas}\label{acciones-permitidas}

\begin{itemize}
\tightlist
\item
  La embarcación tiene una capacidad máxima de dos pasajeros, por lo que
  en cada cruce debe transportar al menos una persona y como máximo dos.
  Los pasajeros pueden ser misioneros, caníbales o una combinación de
  ambos.
\end{itemize}

El conjunto de acciones posibles se representa mediante pares ordenados
\((m,c)\), donde \(m\) indica la cantidad de misioneros y \(c\) la
cantidad de caníbales que participan en el cruce. Por tanto, el conjunto
de acciones es:

\[
A={(1,0),(2,0),(0,1),(0,2),(1,1)}
\]

Cada elemento representa una combinación de pasajeros: un misionero, dos
misioneros, un caníbal, dos caníbales o un misionero junto con un
caníbal, respectivamente.

La dirección del desplazamiento depende de la posición de la
embarcación, representada mediante la variable \(B\) del estado. Cuando
\(B=1\), la embarcación se encuentra en la orilla izquierda y los
pasajeros se trasladan hacia la derecha. Cuando \(B=0\), se encuentra en
la orilla derecha y el desplazamiento se realiza en sentido contrario.

Sin embargo, no todas las acciones pueden ejecutarse desde cualquier
estado. Para que un movimiento sea válido, debe existir en la orilla de
partida una cantidad suficiente de pasajeros para realizarlo. Además, el
estado resultante debe cumplir las restricciones de seguridad
establecidas para el problema.

\subsubsection{Función de transición}\label{funciuxf3n-de-transiciuxf3n}

\begin{itemize}
\tightlist
\item
  La función de transición determina el estado resultante después de
  ejecutar una acción sobre el estado actual. Su comportamiento depende
  de la posición de la embarcación, puesto que esta establece la orilla
  desde la cual se realiza el desplazamiento.
\end{itemize}

Sea el estado actual \(S=(M,C,B)\) y una acción \(a=(m,c)\), donde \(m\)
representa la cantidad de misioneros y \(c\) la cantidad de caníbales
que participan en el cruce. La función de transición se define mediante
dos expresiones, según la ubicación de la embarcación.

Cuando la embarcación se encuentra en la orilla izquierda (\(B=1\)), los
pasajeros se trasladan hacia la derecha y las cantidades de misioneros y
caníbales de la orilla izquierda disminuyen:

\[
T((M,C,1),(m,c))=(M-m,C-c,0)
\]

Por el contrario, cuando la embarcación se encuentra en la orilla
derecha (\(B=0\)), los pasajeros regresan a la izquierda y las
cantidades correspondientes aumentan:

\[
T((M,C,0),(m,c))=(M+m,C+c,1)
\]

En ambos casos, la variable \(B\) cambia de valor para representar el
desplazamiento de la embarcación hacia la orilla opuesta.

La aplicación de esta función está condicionada por dos aspectos:
primero, debe existir una cantidad suficiente de pasajeros en la orilla
de partida para ejecutar la acción; segundo, el estado resultante debe
satisfacer las restricciones de seguridad y los límites establecidos
para las cantidades de misioneros y caníbales. De esta manera,
únicamente se consideran transiciones que conducen a estados válidos
dentro del espacio de búsqueda.

\subsubsection{Restricciones de validez}\label{restricciones-de-validez}

\begin{itemize}
\tightlist
\item
  Para que un estado sea considerado válido, debe cumplir las
  condiciones que garantizan la seguridad de los misioneros en ambas
  orillas del río. La restricción principal establece que los caníbales
  no pueden superar en número a los misioneros en una orilla cuando
  estos últimos se encuentran presentes.
\end{itemize}

Considerando un estado \(S=(M,C,B)\), donde \(M\) y \(C\) representan
las cantidades de misioneros y caníbales en la orilla izquierda, las
cantidades correspondientes a la orilla derecha se obtienen mediante:

\[
M_D=3-M
\]

\[
C_D=3-C
\]

A partir de esta representación, se establecen las siguientes
condiciones:

\begin{enumerate}
\def\labelenumi{\arabic{enumi}.}
\tightlist
\item
  Límites de población: las cantidades de misioneros y caníbales en la
  orilla izquierda deben encontrarse dentro del intervalo de cero a
  tres:
\end{enumerate}

\[
0\leq M\leq 3
\]

\[
0\leq C\leq 3
\]

Además, la variable que representa la posición de la embarcación
únicamente puede tomar los valores cero o uno:

\[
B\in{0,1}
\]

\begin{enumerate}
\def\labelenumi{\arabic{enumi}.}
\setcounter{enumi}{1}
\tightlist
\item
  Seguridad en la orilla izquierda: si existe al menos un misionero en
  esta orilla, su cantidad debe ser igual o superior a la de los
  caníbales:
\end{enumerate}

\[
M=0\quad\text{o}\quad M\geq C
\]

\begin{enumerate}
\def\labelenumi{\arabic{enumi}.}
\setcounter{enumi}{2}
\tightlist
\item
  Seguridad en la orilla derecha: se debe aplicar el mismo criterio en
  la orilla opuesta:
\end{enumerate}

\[
M_D=0\quad\text{o}\quad M_D\geq C_D
\]

Estas condiciones se verifican después de cada transición. Si alguna de
ellas no se cumple, el estado resultante se considera inválido y no debe
incorporarse al espacio de búsqueda.

De esta manera, las restricciones permiten descartar los estados que
representan situaciones prohibidas y garantizan que los algoritmos de
búsqueda exploren únicamente configuraciones válidas del problema.

\subsubsection{Criterios de solución}\label{criterios-de-soluciuxf3n}

\begin{itemize}
\tightlist
\item
  Una solución para el problema de los misioneros y caníbales consiste
  en encontrar una secuencia finita de estados que permita trasladar a
  todos los individuos desde la orilla izquierda hasta la orilla
  derecha, respetando las restricciones establecidas durante cada cruce.
\end{itemize}

La secuencia debe comenzar en el estado inicial \(S_0=(3,3,1)\) y
finalizar en el estado objetivo \(S_G=(0,0,0)\). Por tanto, una solución
puede representarse de la siguiente manera:

\[
S_0\rightarrow S_1\rightarrow S_2\rightarrow\cdots\rightarrow S_k=S_G
\]

Para que esta secuencia sea considerada válida, debe cumplir las
siguientes condiciones:

El primer estado debe corresponder a la configuración inicial del
problema. Cada estado sucesor debe obtenerse mediante la aplicación de
una acción permitida sobre el estado anterior. Todos los estados de la
secuencia deben satisfacer las restricciones de seguridad y los límites
de población. El último estado debe coincidir con el estado objetivo,
indicando que los tres misioneros, los tres caníbales y la embarcación
se encuentran en la orilla derecha.

En este problema, cada cruce de la embarcación tiene el mismo costo,
independientemente de la cantidad y combinación de pasajeros
transportados. Por consiguiente, el costo total de una solución
corresponde al número de transiciones realizadas, representado por
\(k\).

Este criterio permite establecer una diferencia entre los algoritmos de
búsqueda implementados. La búsqueda en anchura (BFS), al explorar los
estados por niveles, garantiza encontrar una solución con el menor
número de cruces cuando todas las acciones tienen el mismo costo. Por
otra parte, la búsqueda en profundidad (DFS) permite encontrar una ruta
válida, pero no garantiza que esta sea la de menor longitud.

De esta manera, el objetivo de la implementación consiste en encontrar
una secuencia de transiciones válidas que alcance el estado objetivo y
analizar el comportamiento de ambos algoritmos durante el proceso de
búsqueda.

\subsection{Implementación en Python}\label{implementaciuxf3n-en-python}

\subsubsection{Implementación de los algoritmos de
búsqueda}\label{implementaciuxf3n-de-los-algoritmos-de-buxfasqueda}

\begin{itemize}
\tightlist
\item
  La formulación del problema como un espacio de estados constituye la
  base para desarrollar su solución computacional. En esta etapa se
  implementaron los elementos definidos anteriormente mediante el
  lenguaje de programación Python, estableciendo mecanismos para
  representar las configuraciones del problema, generar transiciones
  válidas y explorar el espacio de búsqueda.
\end{itemize}

La implementación contempla dos algoritmos de búsqueda ciega: búsqueda
en anchura (BFS) y búsqueda en profundidad (DFS). Ambos utilizan la
misma representación de estados, las mismas acciones permitidas y las
mismas restricciones de seguridad, diferenciándose principalmente en la
estrategia utilizada para explorar los estados disponibles.

\subsubsection{Representación de los estados y
acciones}\label{representaciuxf3n-de-los-estados-y-acciones}

\begin{itemize}
\tightlist
\item
  Cada estado del problema se representa mediante una tupla de tres
  elementos \((M,C,B)\), donde \(M\) y \(C\) indican las cantidades de
  misioneros y caníbales que permanecen en la orilla izquierda, mientras
  que \(B\) identifica la posición de la embarcación. Esta
  representación permite almacenar cada configuración de manera compacta
  y facilita la identificación de estados repetidos durante la búsqueda.
\end{itemize}

El estado inicial se establece mediante la tupla \((3,3,1)\), mientras
que el estado objetivo corresponde a \((0,0,0)\). Asimismo, las
cantidades de individuos en la orilla derecha se calculan a partir de
las cantidades presentes en la izquierda, considerando que la población
total permanece constante.

Para determinar si una configuración puede formar parte del espacio de
búsqueda, se implementó la función estado\_valido(). Esta función
verifica que las cantidades de misioneros y caníbales se encuentren
dentro de los límites establecidos, que la posición de la embarcación
sea válida y que se respeten las restricciones de seguridad en ambas
orillas.

Por otra parte, la función obtener\_sucesores() se encarga de generar
los estados que pueden alcanzarse a partir de una configuración actual.
Para ello, recorre el conjunto de acciones permitidas y comprueba que
exista una cantidad suficiente de pasajeros en la orilla de partida.
Posteriormente, calcula el nuevo estado según la dirección del
desplazamiento y utiliza la función estado\_valido() para determinar si
puede incorporarse como sucesor.

De esta manera, la generación de sucesores evita incorporar
configuraciones que incumplan las condiciones del problema. Ambos
algoritmos utilizan este mismo procedimiento, garantizando que la
comparación de sus estrategias de búsqueda se realice sobre un espacio
de estados común.

\subsubsection{Busqueda en anchura (BFS)}\label{busqueda-en-anchura-bfs}

\begin{itemize}
\tightlist
\item
  La búsqueda en anchura (Breadth-First Search, BFS) se implementó como
  el primer método de exploración del espacio de estados. Este algoritmo
  utiliza una estructura de datos tipo cola para procesar los estados
  según el orden en que fueron descubiertos, siguiendo una estrategia de
  primero en entrar, primero en salir (FIFO).
\end{itemize}

Para su implementación en Python se utilizó la estructura deque,
perteneciente al módulo collections, que permite insertar nuevos estados
al final de la cola y extraerlos desde el inicio de manera eficiente.

El procedimiento comienza incorporando el estado inicial \((3,3,1)\) a
la cola de exploración y al conjunto de estados visitados.
Posteriormente, mientras existan estados pendientes, se extrae el
primero de la cola y se comprueba si corresponde al estado objetivo. Si
no se ha alcanzado la solución, se generan sus sucesores mediante la
función obtener\_sucesores().

Cada sucesor que todavía no ha sido visitado se incorpora al conjunto de
estados visitados y se añade al final de la cola. Adicionalmente, se
registra su estado predecesor mediante una estructura de tipo
diccionario. Este registro permite conservar la relación entre los
estados y reconstruir posteriormente la secuencia de transiciones que
conduce a la solución.

Cuando se encuentra el estado objetivo, el algoritmo utiliza la
información almacenada de los predecesores para reconstruir la ruta
desde el estado final hasta el inicial. Finalmente, invierte la
secuencia obtenida para presentarla en el orden en que deben realizarse
los cruces.

El uso de una cola FIFO garantiza que los estados se exploren por
niveles de profundidad. Por esta razón, dado que todas las transiciones
tienen el mismo costo, BFS permite encontrar una solución con el menor
número de cruces posible.

Si se agotan los estados pendientes sin alcanzar el objetivo, el
algoritmo determina que no se encontró una ruta de solución dentro del
espacio de búsqueda explorado.

\subsection{Busqueda en profundidad
(DFS)}\label{busqueda-en-profundidad-dfs}

La búsqueda en profundidad (Depth-First Search, DFS) se implementó como
el segundo método de exploración del espacio de estados. A diferencia de
BFS, este algoritmo utiliza una estructura de datos tipo pila, siguiendo
una estrategia de último en entrar, primero en salir (LIFO).

En Python, la pila se implementó mediante una lista, utilizando la
operación append() para incorporar nuevos estados y pop() para extraer
el último estado agregado. De esta manera, el algoritmo prioriza la
exploración de los sucesores más recientemente incorporados a la
estructura.

El procedimiento comienza insertando el estado inicial \((3,3,1)\) en la
pila y registrándolo como visitado. Mientras existan estados pendientes
de exploración, se extrae el último elemento de la pila y se comprueba
si corresponde al estado objetivo. Si todavía no se ha alcanzado la
solución, se generan sus sucesores mediante la función
obtener\_sucesores().

Los sucesores que no han sido visitados se incorporan a la pila y se
registran junto con su estado predecesor. Al igual que en BFS, esta
información permite reconstruir la ruta encontrada cuando se alcanza el
estado objetivo. Una vez obtenida, la secuencia se invierte para mostrar
los estados desde la configuración inicial hasta la configuración final.

La estrategia LIFO permite que DFS explore una trayectoria en
profundidad antes de continuar con otras alternativas pendientes. Esto
puede facilitar el descubrimiento de una solución sin necesidad de
explorar previamente todos los estados de niveles anteriores. Sin
embargo, no garantiza encontrar una ruta con el menor número de
transiciones, incluso cuando existe una solución más corta.

En esta implementación, el conjunto de estados visitados evita explorar
repetidamente una misma configuración, mientras que el registro de
predecesores permite recuperar la ruta correspondiente a la primera
solución encontrada. De esta manera, DFS utiliza los mismos mecanismos
de validación y generación de sucesores que BFS, pero modifica el orden
en que se explora el espacio de búsqueda.

\subsection{Metricas de evaluación}\label{metricas-de-evaluaciuxf3n}

Para analizar el comportamiento de los algoritmos de búsqueda
implementados, se establecieron métricas que permiten evaluar el
esfuerzo de exploración, el tamaño de las rutas obtenidas y el tiempo
requerido para encontrar una solución. Estos indicadores facilitan la
comparación entre BFS y DFS, considerando que ambos utilizan el mismo
espacio de estados y las mismas restricciones.

Las métricas consideradas son las siguientes:

\begin{enumerate}
\def\labelenumi{\arabic{enumi}.}
\tightlist
\item
  Nodos expandidos
\end{enumerate}

Representa la cantidad de estados que fueron extraídos de la estructura
de exploración y cuyos sucesores fueron generados y evaluados. Este
indicador permite estimar el trabajo realizado por el algoritmo durante
la búsqueda. No se considera expandido un estado cuando se reconoce como
objetivo antes de generar sus sucesores.

\begin{enumerate}
\def\labelenumi{\arabic{enumi}.}
\setcounter{enumi}{1}
\tightlist
\item
  Estados descubiertos
\end{enumerate}

Corresponde a la cantidad total de estados únicos incorporados a la
estructura de exploración durante la ejecución. Se incluye el estado
inicial y se evita contabilizar nuevamente configuraciones que ya fueron
visitadas. Esta métrica permite observar cuántos estados diferentes
fueron identificados por cada algoritmo.

\begin{enumerate}
\def\labelenumi{\arabic{enumi}.}
\setcounter{enumi}{2}
\tightlist
\item
  Longitud de la solución
\end{enumerate}

Indica la cantidad de transiciones necesarias para alcanzar el estado
objetivo desde el estado inicial. Se obtiene restando uno al número
total de estados que conforman la ruta encontrada:

\[
L=|R|-1
\]

Donde \(R\) representa la secuencia de estados que constituye la
solución. Debido a que cada transición corresponde a un cruce de la
embarcación, esta métrica también representa el número de cruces
realizados.

\begin{enumerate}
\def\labelenumi{\arabic{enumi}.}
\setcounter{enumi}{3}
\tightlist
\item
  Tiempo de ejecución
\end{enumerate}

Mide el tiempo que necesita cada algoritmo para completar la búsqueda y
obtener una solución. Para su medición en Python se utiliza la función
perf\_counter() del módulo time, que permite registrar intervalos de
tiempo con alta resolución.

El tiempo de ejecución se calcula mediante la diferencia entre los
instantes registrados al finalizar y comenzar el algoritmo:

\[
T_e=t_{\text{final}}-t_{\text{inicial}}
\]

Este indicador permite comparar el tiempo requerido por cada estrategia
de búsqueda bajo las mismas condiciones de ejecución.

\begin{enumerate}
\def\labelenumi{\arabic{enumi}.}
\setcounter{enumi}{4}
\tightlist
\item
  Factor de ramificación efectivo
\end{enumerate}

El factor de ramificación efectivo (\(b^*\)) es un indicador que permite
representar el crecimiento de los nodos descubiertos en relación con la
profundidad de la solución. Se obtiene mediante la siguiente expresión:

\[
N = 1 + b^{*} + (b^{*})^2 + \cdots + (b^{*})^d
\]

Donde \(N\) representa el número total de estados descubiertos,
incluyendo el estado inicial, y \(d\) corresponde a la profundidad de la
solución, equivalente a la cantidad de transiciones realizadas.

El valor de \(b^*\) se determina mediante una aproximación numérica que
satisface la igualdad anterior. Este indicador permite describir el
esfuerzo de exploración en función de la profundidad alcanzada; sin
embargo, debe interpretarse considerando que se trata de una estimación
y que el espacio de estados del problema es finito.

En conjunto, estas métricas permiten evaluar tanto la eficiencia de
exploración como la calidad de las soluciones obtenidas. Su utilización
proporciona una base cuantitativa para analizar las diferencias entre
BFS y DFS durante la resolución del problema de los misioneros y
caníbales.

\subsection{Resultados experimentales}\label{resultados-experimentales}

Una vez implementados los algoritmos de búsqueda en anchura (BFS) y
búsqueda en profundidad (DFS), se realizaron las ejecuciones
correspondientes utilizando la misma configuración del problema, el
mismo estado inicial y las mismas condiciones de transición.

Los resultados obtenidos permiten analizar las rutas encontradas y
comparar el comportamiento de ambos algoritmos mediante las métricas
establecidas. Para ello, se registraron los estados recorridos, los
nodos expandidos, los estados descubiertos, la longitud de las
soluciones, el tiempo de ejecución y el factor de ramificación efectivo.

Rutas encontradas por los algoritmos

Ambos algoritmos consiguieron encontrar una secuencia de transiciones
válidas que conecta el estado inicial \((3,3,1)\) con el estado objetivo
\((0,0,0)\). Sin embargo, las rutas obtenidas presentan diferencias en
algunas de las configuraciones intermedias debido al orden de
exploración utilizado por cada estrategia.

La siguiente tabla muestra los estados que conforman las soluciones
encontradas. Cada fila representa una etapa del recorrido, donde el
número de cruce indica la cantidad de transiciones realizadas desde el
estado inicial.

\subsubsection{Tabla de rutas
encontradas}\label{tabla-de-rutas-encontradas}

{\def\LTcaptype{none} % do not increment counter
\begin{longtable}[]{@{}lll@{}}
\toprule\noalign{}
Cruce & BFS & DFS \\
\midrule\noalign{}
\endhead
\bottomrule\noalign{}
\endlastfoot
0 & (3,3,1) & (3,3,1) \\
1 & (3,1,0) & (2,2,0) \\
2 & (3,2,1) & (3,2,1) \\
3 & (3,0,0) & (3,0,0) \\
4 & (3,1,1) & (3,1,1) \\
5 & (1,1,0) & (1,1,0) \\
6 & (2,2,1) & (2,2,1) \\
7 & (0,2,0) & (0,2,0) \\
8 & (0,3,1) & (0,3,1) \\
9 & (0,1,0) & (0,1,0) \\
10 & (1,1,1) & (0,2,1) \\
11 & (0,0,0) & (0,0,0) \\
\end{longtable}
}

Se observa que ambas rutas alcanzan el estado objetivo después de
realizar 11 cruces. Aunque existen diferencias en algunas
configuraciones intermedias, las dos secuencias cumplen las
restricciones de seguridad y las condiciones establecidas para el
desplazamiento de la embarcación.

En el caso de BFS, la solución encontrada corresponde a una ruta de
longitud mínima, debido a que el algoritmo explora los estados por
niveles y todas las transiciones tienen el mismo costo. Por su parte,
DFS también encontró una solución de igual longitud en esta ejecución
particular; sin embargo, este resultado no representa una garantía del
algoritmo, ya que su estrategia de exploración no asegura encontrar
rutas óptimas.

\subsection{Comparación de
métricas}\label{comparaciuxf3n-de-muxe9tricas}

Para evaluar el comportamiento de ambos algoritmos se recopilaron los
valores obtenidos durante la ejecución del programa. La comparación
considera tanto el esfuerzo de exploración como la longitud de las rutas
y el tiempo requerido para encontrar una solución.

{\def\LTcaptype{none} % do not increment counter
\begin{longtable}[]{@{}lll@{}}
\toprule\noalign{}
Métrica & BFS & DFS \\
\midrule\noalign{}
\endhead
\bottomrule\noalign{}
\endlastfoot
¿Existe ruta? & Sí & Sí \\
Nodos expandidos & 14 & 11 \\
Estados descubiertos & 15 & 15 \\
Longitud de la solución & 11 & 11 \\
Tiempo de ejecución (segundos) & 0.000230800 & 0.000066200 \\
Factor de ramificación efectivo & 1.039700 & 1.039700 \\
\end{longtable}
}

Los resultados muestran que ambos algoritmos consiguieron resolver el
problema y descubrieron la misma cantidad de estados. No obstante, se
identifican diferencias en el número de nodos expandidos: BFS expandió
14 estados, mientras que DFS expandió 11. Esto indica que, durante esta
ejecución, DFS necesitó examinar una menor cantidad de estados antes de
encontrar el objetivo.

Respecto a la longitud de las soluciones, ambos algoritmos obtuvieron
rutas de 11 transiciones. Este resultado permite observar que DFS
encontró, en esta ejecución, una solución con la misma longitud que la
obtenida mediante BFS, aunque no posee una garantía general de
optimalidad.

En cuanto al tiempo de ejecución, BFS registró aproximadamente
\(0.000231\) segundos, mientras que DFS registró \(0.0000662\) segundos.
Por tanto, en esta ejecución DFS presentó un menor tiempo de
procesamiento. Sin embargo, debido a que los tiempos registrados son muy
pequeños y corresponden a una ejecución individual, esta diferencia debe
interpretarse como un resultado puntual y no como una conclusión
definitiva sobre el rendimiento general de ambos algoritmos.

Finalmente, el factor de ramificación efectivo presentó el mismo valor
para ambos métodos, equivalente a \(1.039700\). Esto se debe a que su
cálculo utiliza la cantidad de estados descubiertos y la profundidad de
la solución, valores que fueron iguales en las dos ejecuciones.

\subsection{Conclusiones}\label{conclusiones}

\begin{itemize}
\item
  La formulación del problema de los misioneros y caníbales como un
  espacio de estados permitió representar de manera estructurada las
  diferentes configuraciones posibles y establecer las condiciones
  necesarias para alcanzar una solución. La definición del estado, las
  acciones permitidas, la función de transición y las restricciones de
  seguridad constituyeron los elementos fundamentales para desarrollar
  la implementación computacional.
\item
  La implementación de los algoritmos de búsqueda en anchura (BFS) y
  búsqueda en profundidad (DFS) permitió comprobar cómo dos estrategias
  de búsqueda ciega pueden explorar un mismo espacio de estados
  siguiendo órdenes diferentes. Ambos algoritmos consiguieron encontrar
  una solución válida, alcanzando el estado objetivo mediante 11
  transiciones y descubriendo 15 estados únicos durante sus respectivas
  ejecuciones.
\item
  La comparación de las métricas evidenció diferencias en el esfuerzo de
  exploración. BFS expandió 14 nodos, mientras que DFS expandió 11, y
  este último registró un menor tiempo de ejecución en el experimento
  realizado. Sin embargo, estos resultados corresponden a una ejecución
  particular y no permiten establecer que DFS sea generalmente más
  eficiente. Por otra parte, aunque ambos algoritmos encontraron
  soluciones con la misma longitud, BFS garantiza una solución de
  longitud mínima cuando todas las transiciones tienen el mismo costo,
  mientras que DFS no ofrece esta garantía.
\end{itemize}




\end{document}
